\subsection*{What the dictionary adds over the principal subspace}
\label{sec:svd}

A dictionary of 4,608 atoms and a basis of 50 principal axes are not comparable objects, so the
observation that few features align with the leading axes cannot by itself establish that the
information is inaccessible to linear decomposition. We therefore computed the exact
eigendecomposition of the activation covariance at every layer and compared the two representations
under matched conditions (Fig.~\ref{fig:svd}, \suptab{svd_sweep} and \suptab{svd_capacity}).

\paragraph{Reconstruction-matched capacity.} Truncated SVD requires rank \textbf{250--403} to reach
the variance the SAE explains on held-out positions---403 at layer~0 (SAE 0.848), 317 at layer~5
(0.856), 308 at layer~11 (0.823) and 250 at layer~17 (0.775). A 50-axis basis is five- to eightfold
short of the SAE's own reconstruction capacity, so it is the wrong reference for asking what the
dictionary can see.

\paragraph{Where SAE directions sit.} For each unit-normalized decoder direction we computed the
cumulative projection $\rho_k$ onto the rank-$k$ principal subspace, for every $k$ up to the full
1,152 dimensions, against the random-direction expectation $\rho_k = k/d$. SAE directions are
systematically concentrated in the leading subspace: at $k{=}100$ the median direction has
$\rho_k = 0.21$ at layer~0 and 0.33 at layer~11, against 0.087 for a random direction, a two- to
fourfold excess. The median direction reaches half of its norm by $k = 247$ (L0), 190 (L5), 178
(L11) and 152 (L17), where a random direction requires 576, and reaches 90\% by $k = 664$--764
against 1,037. SAE features are not hidden from the principal subspace; they occupy its
high-variance part preferentially.

\paragraph{Low single-axis alignment is not a capacity artefact.} The fraction of features whose
absolute cosine with at least one of the leading $k$ axes exceeds 0.5 saturates by $k \approx
20$--50 and does not increase thereafter. Even against the \emph{complete} 1,152-dimensional basis
it is 0.98\% at layer~0, 0.13\% at layer~5, 0.09\% at layer~11 and 0.26\% at layer~17. Offering more
axes does not recruit more features, because each feature distributes its norm across many axes
rather than concentrating on one. The same conclusion holds across cosine thresholds
$\{0.3, 0.5, 0.7, 0.9\}$ (\suptab{svd_sweep}).

\paragraph{Annotation yield per direction.} We annotated the leading 50 principal axes with the
identical top-20-gene enrichment protocol used for SAE features, in both polarities, since a
principal axis is signed whereas a TopK feature is one-sided. Per direction, principal axes are at
least as annotatable as features: at layer~0, 98\% of axes carry at least one significant term with
a mean of 159.5 terms per direction and 2,029 unique terms, against 90\%, 64.4 and 1,977 for 100
randomly drawn SAE features; at layer~11 the corresponding numbers are 98\%, 178.8 and 2,083 against
93\%, 71.3 and 2,045. Biological annotation is therefore not the exclusive property of
non-aligned directions.

Taken together, these three measurements locate what the dictionary contributes. It is not access to
information that linear decomposition cannot reach---the leading principal subspace contains most of
each feature's norm and annotates readily. It is a different \emph{parameterization} of that
information: thousands of sparse, one-sided, individually interpretable directions in place of a few
hundred dense, signed, polysemantic ones. That is the property the rest of this study relies on,
because it is what makes per-feature annotation, per-feature co-activation structure, and
per-feature causal intervention possible at all.

\subsection*{How many distinct concepts the atlas contains}
\label{sec:concepts}

Dictionary atoms trained at different layers occupy different vector spaces, so their counts cannot
be pooled into a number of concepts held simultaneously in one representation. What can be compared
across layers is each feature's gene-level signature. We built a similarity graph over all 82,525
Geneformer features from their top-20 gene sets and partitioned it (Methods; Fig.~\ref{fig:concepts},
\suptab{concepts}).

The 82,525 atoms resolve into \textbf{18,711 distinct gene-level programs}, that is 4.4 atoms per
program. The count is insensitive to the clustering resolution: 18,695 at $\gamma{=}0.5$ and 18,748
at $\gamma{=}2.0$. The distribution is strongly bimodal: 17,690 programs are singletons---features
whose gene signature overlaps no other feature's---while 1,021 multi-atom programs absorb the
remaining 64,835 atoms, the largest containing 3,221. Redundancy is therefore concentrated in a
small number of heavily re-derived programs rather than spread evenly. As exact sets the signatures
are nearly all distinct: 82,442 distinct signatures, with only 56 shared by more than one feature.

Programs recur across depth. Of the multi-atom programs, 97.1\% appear at two or more layers, and
the mean program spans 3.16 layers. Read together with the near-complete representational turnover
between adjacent layers, this gives a consistent picture: each layer rebuilds a largely fresh
dictionary, but the gene-level programs those dictionaries express are re-derived throughout the
stack.

Within a single layer---the only place a features-per-space ratio is defined---the dictionary is
4$\times$ overcomplete (4,608 atoms in 1,152 dimensions, of which 4,598 are alive at layer~11) and
resolves 774--2,048 distinct programs depending on depth (1,301 at L0, 814 at L5, 1,252 at L11,
1,360 at L17). The packing is close to the geometric limit: mean absolute decoder coherence is
0.036 at layer~11 against a Welch bound of 0.0255 for 4,608 unit vectors in 1,152 dimensions, so
the dictionary is near-optimally spread rather than exploiting an unusually large capacity.
